\subsubsection{Описание и кодирование входных воздействий}
\label{sec:inputs}

Входы $X = \{a,b,c,\tau,\rho,\varnothing\}$ кодируются тремя разрядами $x_2x_1x_0$ (таблица~\ref{tab:inputs}).

\begin{table}[H]
	\centering
	\caption{Кодирование входных воздействий}
	\label{tab:inputs}
	\small\setstretch{1.12}
	\renewcommand{\arraystretch}{1.12}
	\begin{tabularx}{\textwidth}{|c|>{\raggedright\arraybackslash}X|c|}
		\hline
		Вход & Воздействие & Код \\
		\hline
		$a$ & Нажатие кнопки $a$ & $000$ \\
		$b$ & Нажатие кнопки $b$ & $001$ \\
		$c$ & Нажатие кнопки $c$ & $010$ \\
		$\tau$ & Начало четверти часа (минуты $00$, $15$, $30$, $45$) & $100$ \\
		$\rho$ & Истечение одной секунды после входа в $S_{12}$ & $101$ \\
		$\varnothing$ & Отсутствие события & $110$ \\
		\hline
	\end{tabularx}
\end{table}

Событие $\tau$ формируется при естественном ходе часов в начале очередной четверти часа при нулевом значении секунд. Событие $\rho$ вырабатывает таймер, запущенный при входе в $S_{12}$.
